\chapter{Zusammenfassung und Fazit}
\label{cha:Zusammenfassung}




%
%Auf zwei bis drei Seiten soll auf folgende Punkte eingegangen werden:
%
%\begin{itemize}
%	\item Welches Ziel sollte erreicht werden
%	\item Welches Vorgehen wurde gewählt
%	\item Was wurde erreicht, zentrale Ergebnisse nennen, am besten quantitative Angaben machen
%	\item Konnten die Ergebnisse nach kritischer Bewertung zum Erreichen des Ziels oder zur Problemlösung beitragen
%	\item  Ausblick
%\end{itemize}
%
%In der Zusammenfassung sind unbedingt klare Aussagen zum Ergebnis der Arbeit zu nennen. Üblicherweise können Ergebnisse nicht nur qualitativ, sondern auch quantitativ benannt werden, z.~B. \glqq \ldots konnte eine Effizienzsteigerung von \SI{12}{\percent} erreicht werden.\grqq~oder \glqq \ldots konnte die Prüfdauer um \SI{2}{\hour} verkürzt werden\grqq.
%
%Die Ergebnisse in der Zusammenfassung sollten selbstverständlich einen Bezug zu den in der Einleitung aufgeführten Fragestellungen und Zielen haben.

Ziel dieser Studienarbeit war die Weiterentwicklung eines Demonstrationsroboters für den Einsatz auf Hochschulmessen. Aufbauend auf dem in einer vorherigen Studienarbeit entwickelten Funktionsdemonstrator wurden sowohl die mechanische Konstruktion als auch das äußere Erscheinungsbild, die Sensorik sowie die Interaktionsmöglichkeiten des Roboters erweitert und verbessert.

Im mechanischen Bereich wurde die Konstruktion grundlegend überarbeitet. Durch die Neupositionierung des Stützrades konnte die Standsicherheit des Roboters verbessert werden. Zusätzlich wurden ein vollständig neues Verkleidunge für das Gehäuse, ein interaktiver Roboterkopf sowie bewegliche Roboterarme entwickelt und mithilfe additiver Fertigung realisiert. Die modulare Konstruktion ermöglicht dabei zukünftige Erweiterungen und vereinfacht Wartungs- sowie Anpassungsarbeiten.

Zur Erhöhung der Funktionalität wurden weitere Komponenten in das Gesamtsystem integriert. Hierzu zählen insbesondere Ultraschallsensoren zur Hinderniserkennung sowie Servomotoren für die Schulterbewegungen und die Robotergreifer. Die Auswahl der jeweiligen Konzepte erfolgte anhand mehrerer Trade-off-Analysen, wodurch die getroffenen Konstruktionsentscheidungen nachvollziehbar begründet werden konnten.

Ein weiterer Schwerpunkt der Arbeit lag auf der Verbesserung der Mensch-Roboter-Interaktion. Durch die Integration eines sprachbasierten Dialogsystems auf Basis eines Large Language Models sowie eines Displays zur Darstellung animierter Augen wurde eine deutlich natürlichere Kommunikation mit Messebesuchern ermöglicht. Gleichzeitig wurden die Grenzen aktueller Sprachmodelle, insbesondere hinsichtlich Antwortlatenz und Halluzinationen, untersucht und bewertet.

Die beiden Messeeinsätze auf der Explore Science Friedrichshafen sowie dem Technikforum der DHBW Friedrichshafen ermöglichten erstmals eine Erprobung des Roboters unter realen Einsatzbedingungen. Dabei zeigte sich, dass das entwickelte Designkonzept von den Besuchern sehr positiv aufgenommen wurde. Insbesondere Kinder zeigten großes Interesse am Roboter und suchten aktiv die Interaktion. Gleichzeitig offenbarten die Praxiseinsätze mehrere technische Schwachstellen, unter anderem in der Spannungsversorgung, der Zuverlässigkeit einzelner Hardwarekomponenten sowie der Sprachausgabe. Diese Erkenntnisse stellen einen wichtigen Beitrag für die zukünftige Weiterentwicklung des Systems dar.

Insgesamt konnten die im Rahmen dieser Studienarbeit gesetzten Ziele weitgehend erreicht werden. Das Erscheinungsbild des Roboters wurde deutlich verbessert und durch zusätzliche Interaktionsmöglichkeiten erweitert. Gleichzeitig konnte gezeigt werden, dass sich der entwickelte Demonstrator grundsätzlich für den Einsatz auf Messen eignet. Die praktischen Erprobungen machten jedoch deutlich, dass für einen dauerhaft zuverlässigen Betrieb insbesondere die Robustheit der Hardware sowie die Stabilität der Software weiter verbessert werden müssen.

\section*{Ausblick}

Die im Rahmen dieser Studienarbeit gewonnenen Erkenntnisse bilden die Grundlage für weitere Entwicklungsarbeiten. Ein Schwerpunkt zukünftiger Arbeiten sollte auf der Erhöhung der Betriebssicherheit und Zuverlässigkeit des Gesamtsystems liegen. Hierzu zählen insbesondere eine robustere Auslegung der Spannungsversorgung, eine verbesserte Absicherung der Elektronik sowie eine umfassende Fehlerdiagnose zur schnelleren Identifikation von Hardwareproblemen.

Darüber hinaus bietet das entwickelte Robotersystem Potenzial für zahlreiche funktionale Erweiterungen. Hierzu gehören beispielsweise eine autonome Fahrfunktion, die mit Ultraschallsensoren und kamerabasierte Objekt- und Personenerkennung arbeitet, sowie die Integration funktionierender Gesten und Bewegungsabläufe. Auch eine Optimierung des Sprachdialogsystems, etwa durch leistungsfähigere Sprachmodelle.

Durch die modulare Konstruktion und den hohen Anteil additiv gefertigter Komponenten können zukünftige Erweiterungen mit vergleichsweise geringem Aufwand umgesetzt werden. Der entwickelte Roboter stellt somit eine geeignete Plattform für weitere Studienarbeiten dar und kann kontinuierlich an neue technische Möglichkeiten und Anforderungen angepasst werden.